% title:          Exponential equation with three solutions
% area:           real-analysis, algebraic-identities
% methods:        rolle-theorem, contradiction
% difficulty:
% difficulty-ai:  medium
% status:
% review:         none
% --- history: all optional, leave EMPTY if unknown ---
% origin:
% origin-date:
% origin-number:
% also-in:
% taken-from:
% references:     Method background (the rule of signs for exponential sums, which the solution's Rolle argument re-proves): G.J.O. Jameson, 'Counting zeros of generalized polynomials: Descartes' rule of signs and Laguerre's extensions', Math. Gazette 90 (2006) 223-234 - https://www.maths.lancs.ac.uk/~jameson/zeros.pdf; related problem (not the same equation): Putnam and Beyond (1st ed.), Problem 423, 6^x+1=8^x-27^(x-1) - https://mathematicalolympiads.wordpress.com/wp-content/uploads/2012/08/putnam-and-beyond.pdf, solved with Laguerre's rule at Math StackExchange 800319 - https://math.stackexchange.com/questions/800319
% history:        ai
\begin{problem}
Find all real solutions to the equation
\[
9^{x} + 4^{x}+2^{x} = 8^{x} + 6^{x}+1.
\]
\end{problem}
\begin{solution}
It is easy to see that \( x = 0 \), \( x = 1 \), and \( x = 2 \) are solutions. So the equation has at least 3 distinct real solutions. Let us introduce the function \( f: \mathbb{R} \rightarrow \mathbb{R} \) defined by
\[
f(x) = 9^{x} + 4^{x} + 2^{x} - 8^{x} + 6^{x} + 1.
\]
As stated above, \( f \) has at least 3 distinct zeros. We claim there are no other roots. By Rolle's theorem, if a function \( g: \mathbb{R} \rightarrow \mathbb{R} \) has at least \( n \geq 2 \) zeros \( x_1 < \dots < x_n \), then the function \( g'(x) \) has at least \( n - 1 \) zeros \( y_1 < \dots < y_{n-1} \), where for each \( i = 1, \dots, n-1 \), we have \( x_i < y_i < x_{i+1} \). In particular, since for each \( a \in \mathbb{R}_{>0} \), \( a^{-x} g(x) \) has at least \( n \) zeros, we have that \( \left( a^{-x} g(x) \right)' \) has at least \( n - 1 \) zeros, and so does the function
\[
h_{a} g(x) = a^x \left( a^{-x} g(x) \right)' = g'(x) - \ln(a) g(x).
\]
Suppose \( f \) has another zero not in \( \{0, 1, 2\} \). Then \( f \) has at least \( 4 \) zeros, and thus
\[
h_{1} f(x) = f'(x) = \ln(9) 9^x + \ln(4) 4^x + \ln(2) 2^x - \ln(8) 8^x - \ln(6) 6^x
\]
has at least \( 3 \) zeros, which then implies that
\[
h_{6} h_{1} f(x) = f''(x) - \ln(6) f'(x)
\]
\[
= \ln\left( \frac{9}{6} \right) \ln(9) 9^x + \ln\left( \frac{4}{6} \right) \ln(4) 4^x + \ln\left( \frac{2}{6} \right) \ln(2) 2^x - \ln\left( \frac{8}{6} \right) \ln(8) 8^x
\]
has at least \( 2 \) zeros, which again implies that
\[
h_{8} h_{6} h_{1} f(x) = \ln\left( \frac{9}{8} \right) \ln\left( \frac{9}{6} \right) \ln(9) 9^x - \ln\left( \frac{4}{8} \right) \ln\left( \frac{4}{6} \right) \ln(4) 4^x - \ln\left( \frac{2}{8} \right) \ln\left( \frac{2}{6} \right) \ln(2) 2^x
\]
has at least \( 1 \) zero.
\\\\
The function \( h_{8} h_{6} h_{1} f(x) \) is of the form \( k_{2} 2^x + k_{4} 4^x + k_{9} 9^x \), for $k_2, k_4, k_9>0$ and hence is always positive. Therefore, \( h_{8} h_{6} h_{1} f(x) \) cannot have any real zero. This is a contradiction to the assumptions hence the solutions of the original equation are exactly \( \{0, 1, 2\} \).
\end{solution}
